\documentclass{report}
\usepackage{amsmath,amssymb}
\setlength{\parindent}{0mm}
\setlength{\parskip}{1em}
\begin{document}
\begin{center}
\rule{6in}{1pt} \
{\large Johnny Valbuena \\
{\bf The shape grain influence on the permittivity: Numerical study}}

Department of Applied Mathemactics (RSPE) \\ Australian National University \\ L 1 24 Le Couteur Building \\ Canberra 200 \\ Australia
\\
{\tt johnny.valbuena@anu.edu.au}\end{center}

Electric phenomena occur as a response of the dielectric material under
the influence of an electric field at different frequencies such as
polarization and relaxation processes. They establish the dielectric
properties which can be used to characterize the material properties, for
instance moisture content, bulk density, bio-content and chemical
concentration. The relationship between them plays an important role for
research and application in food science, medicine, biology, agriculture,
chemistry, electric devices, defence industry, engineering. One important
reason for the interest in dielectric responses of rocks lies in the
investigation of their physical properties in a non-destructive manner at
a considerable lower cost. They are used in the petroleum industry to
estimate reservoir parameters which are important to study the reservoir
formation, evaluate zones for hydrocarbon reserves and oil recovery
projects.

Permittivity is a property of the dielectric material which measures the
ability of the material to be polarized by an electric field. In a static
state, it is defined as ${\bf D} = \varepsilon {\bf E}$ where
$\varepsilon$ is the permittivity, ${\bf D}$ and ${\bf E}$ are the
electric flux density and electric field, respectively. This equation
holds for linear, homogeneous and isotropic materials. For anisotropic
materials, the permittivity becomes a second rank tensor. When the
material consists of dielectric and conductivity compounds, and an
alternating electric field is applied, ${\bf D}$ and ${\bf E}$ are not in
phase. Then, the permittivity is defined as $ \varepsilon^* = \big (
\varepsilon_r \varepsilon_o - j \frac{\sigma}{ \omega} \big )$ and it is
called complex permittivity. There are several ways to calculate the
effective permittivity of the medium, which are: total current ${\bf J}$
and the phase difference $\theta$, Gauss' law, energy balance and using
average values of the electric displacement $<{\bf D}>$ and the electric
field $<{\bf E}>$. Also, it is necessary to calculate the distribution of
the potentials in the medium by using the continuity equation for the
current density $\big ( \nabla \cdot \big [ \big ( \varepsilon_r
\varepsilon_o \big ) \nabla \Phi \big ] = 0 \big ) $.

This research is focused on the influence of the shape grain on the
permittivity at different frequencies using 3D granular models, 3D images
of porous materials, mixing law and Finite Element Method. The finite
element is represented by a voxel with 8 nodes, one in each corner. In
each node an electrical potential is applied, then the approximation of
the potential ($\phi_e$) within an element is determined by the
tri--linear interpolation and interrelates the potential distribution in
various elements such that the potential is continuous across
interelement boundaries. The interpolation scheme involves 26 neighbours
and the interpolated potential is expressed as $\phi_{e}(x,y,z) =
\sum_{i=1}^{8} \alpha_{i}(x,y,z) \phi_i $, where $\alpha$ is the
interpolation function. The electric field in the voxel is obtained by
${\bf E}_e = -\nabla \phi_{e}(x,y,z)$. The function of energy
corresponding to the equation of current density is $W_{e} = \frac{1}{2}
\int_{vol} \epsilon | {\bf E_e} | ^2 d(vol) $. When the process of
assembling over all elements of the material is carried out, the total
energy is given by ${\bf W} = \sum_{e=1}^{N} = {\bf W_e} = \frac{1}{2}
\epsilon [\Phi]^T [C] [\Phi] $, where $\Phi$ is a vector and $C$ the
global stiffness matrix. The current density equation is satisfied when
the total energy in the solution is minimum, then it requires that the
partial derivative of ${\bf W}$ with respect to each node value of the
potential be zero $ \big ( \frac{\partial {\bf W}}{ \partial \phi_k} =
{\bf 0} \big ) $.

A system of equations $ A \Phi = {\bf b} $ is generated by $
\frac{\partial {\bf W}}{ \partial \phi_k} = {\bf 0} $, where $A$ is a
sparse matrix which represents the global stiffness matrix, $\Phi$ is a
vector with all the potentials whose components depend on three
coordinates in the image as well as vector ${\bf b}$, which denotes the
boundary conditions. This system is solved in order to minimise the
potential and to calculate the total energy. Dirichlet boundary
conditions are used on the bottom and top of the 3D image and Neumann
boundary conditions on the other faces of the image. The boundary
conditions are represented by voltage which causes an electric field
across of the image. When an static field is applied, the matrix $A$ is
symmetric owing to Laplacian operator of the current density equation.
However, a system of complex equations is generated when an alternating
electric field is applied, thus the matrix $A$ is not Hermitian but
symmetric.

The general procedure to calculate an effective permittivity from a image
is as follows: apply the voltage, calculate the local and global
stiffness matrices, solve the equations system and then the property can
be calculated by methods mentioned above. The dielectric constant of the
material within each voxel is known and local potentials are already
calculated after solving the system of equations. Then, the local
electric field and local electric flux density can be calculated as well.
The last step was to use the average values of $<{\bf D}>$ and $<{\bf
E}>$ where the effective permittivity is given by $\varepsilon_{eff} =
<{\bf D}>/<{\bf E}> $. Following this procedure, the numerical results
fit well with mixing laws for samples (a cube with a sphere at its
center) in different sizes to 80 Voxels and in a static field. The system
of equations was solved using the algorithms BICG, GMRES, QMR and TFQMR.
We are working on using these algorithms to solve a system of complex
equations which represents the main difficulty. We need to utilize
precondition and domain decomposition techniques in order to increase the
size of the image that is usually between 2000 and 3000 voxels. Thus,
research of the effective permittivity of the material would become more
useful and interesting.


\end{document}
